\documentclass[11pt]{article}
\usepackage[utf8]{inputenc}
\usepackage[margin=1in]{geometry}
\usepackage{amsmath}
\usepackage{graphicx}
\usepackage{booktabs}
\usepackage{hyperref}
\usepackage[round]{natbib}
\usepackage{float}
\usepackage{multirow}
\usepackage{array}
\usepackage{siunitx}

\title{Shape-Changing Electrode Array for Minimally Invasive Large-Scale Intracranial Brain Activity Mapping}

\author{%
  Author A$^{1,*}$, Author B$^{1}$, Author C$^{2}$, Author D$^{3}$ \\[6pt]
  $^{1}$Department of Electrical and Computer Engineering, Rice University, Houston, TX, USA \\
  $^{2}$Department of Bioengineering, Rice University, Houston, TX, USA \\
  $^{3}$Department of Neurosurgery, Baylor College of Medicine, Houston, TX, USA \\
  $^{*}$Corresponding author: author.a@rice.edu
}

\date{}

\begin{document}
\maketitle

\begin{abstract}
Large-scale brain activity mapping with high spatiotemporal resolution requires electrocorticography (ECoG), but current implantation methods demand extensive craniotomy and durotomy with significant surgical risks. We developed a shape-changing electrode array (SCEA) that transforms from a compressed 3-mm-wide strip into a $20 \times 15$~mm$^2$ sheet upon reaching body temperature (37$^\circ$C), enabling minimally invasive implantation through small skull openings. The SCEA uses a CNT/Au multilayer structure on a 4-$\mu$m Parylene-C substrate with a nitinol shape-memory alloy actuator. The 64-channel array maintained stable impedance before and after shape transformation (pre: $48.7 \pm 6.2$~k$\Omega$; post: $49.3 \pm 5.8$~k$\Omega$ at 1~kHz; paired $t$-test $p = 0.42$, $n = 64$ channels). Longitudinal MRI in rats showed no brain structural changes at any timepoint (1--8 weeks), and meningeal enhancement on DCE-MRI resolved completely by 4 weeks. Immunohistochemistry confirmed minimal GFAP ($1.12 \pm 0.08$ fold over contralateral at 8 weeks) and Iba1 ($1.09 \pm 0.06$ fold) responses. The SCEA captured high-frequency oscillations up to 500~Hz during 4-aminopyridine-induced seizures in rats and recorded large-scale cortical dynamics during anesthesia emergence in a 256-channel canine subdural configuration. This technology enables large-scale, high-bandwidth brain mapping with substantially reduced surgical invasiveness.

\medskip
\noindent\textbf{Keywords:} electrocorticography, shape-memory alloy, minimally invasive, neural electrode, brain-computer interface, carbon nanotube, high-frequency oscillations
\end{abstract}

\section{Introduction}

Understanding brain function at the scale of cortical networks requires recording techniques with high spatial resolution (sub-millimeter), high temporal resolution (sub-millisecond), and broad bandwidth ($>$500~Hz) \citep{buzsaki2012origin, engel2005invasive}. Electrocorticography (ECoG) uniquely satisfies these requirements by placing electrodes directly on the cortical surface, enabling detection of high-frequency oscillations (HFOs), local field potentials, and even single-unit activity \citep{khodagholy2015neurogrid, viventi2011flexible}. ECoG has proven indispensable for epilepsy surgery planning \citep{litt2001epileptic, worrell2004high}, brain-computer interfaces \citep{schalk2007decoding}, and fundamental neuroscience research.

However, conventional ECoG implantation requires extensive craniotomy and durotomy---surgical procedures that carry significant risks including infection, inflammation, edema, and hemorrhage \citep{lesser2010subdural}. These risks scale with the size of the electrode array, creating a fundamental tension: larger arrays provide more comprehensive brain coverage but demand more invasive surgery. This limits the clinical translation of large-scale ECoG for conditions ranging from epilepsy to neural prosthetics.

Noninvasive alternatives cannot match ECoG's capabilities. Scalp EEG provides centimeter-scale spatial resolution with limited bandwidth ($<$50~Hz in routine clinical use) \citep{buzsaki2012origin}. Functional MRI and fNIR offer millimeter spatial resolution but are limited to hemodynamic signals with temporal resolution on the order of seconds \citep{fox2009column}. None of these modalities can detect the HFOs (80--500~Hz) that are critical biomarkers for epileptogenic zones \citep{jacobs2009high}.

Recent advances in flexible electronics \citep{rogers2010materials, kim2010dissolvable} and novel electrode materials \citep{wang2019graphene, cogan2008neural} have improved the conformability and biocompatibility of neural interfaces. Carbon nanotube (CNT) electrodes maintain conductivity under mechanical deformation \citep{elgamal2014carbon}, and ultrathin polymer substrates such as Parylene-C provide biocompatible, flexible platforms \citep{luo2023parylene}. Shape-memory alloys (SMAs), particularly nitinol (NiTi), undergo temperature-driven phase transitions that can actuate mechanical shape changes \citep{jang2019nitinol}.

Here, we combine these technologies into a shape-changing electrode array (SCEA) that can be inserted through a small slit in the skull or dura in a compressed strip configuration and then self-deploys into a large-area sheet at body temperature. We demonstrate stable electrical performance, biocompatibility through longitudinal MRI and immunohistochemistry, high-bandwidth seizure recording in rats, and large-scale cortical mapping in a canine model.

\section{Results}

\subsection{SCEA Design and Specifications}

Table~\ref{tab:design} summarizes the SCEA design parameters.

\begin{table}[H]
\centering
\caption{SCEA design specifications for rat (64-channel) and canine (256-channel) configurations. CNT = carbon nanotube; SWCNT = single-walled CNT; CVD = chemical vapor deposition; PEO = polyethylene oxide.}
\label{tab:design}
\begin{tabular}{lcc}
\toprule
\textbf{Parameter} & \textbf{Rat (64-ch)} & \textbf{Canine (256-ch)} \\
\midrule
Substrate material & Parylene-C & Parylene-C \\
Substrate thickness ($\mu$m) & 4 & 4 \\
Insulation layers & SU8 & SU8 \\
Conductive layer & CNT/Au & CNT/Au \\
Active site size ($\mu$m$^2$) & $100 \times 100$ & $100 \times 100$ \\
Compressed strip width (mm) & 3 & 3 \\
Deployed array size (mm$^2$) & $20 \times 15$ & $20 \times 10$ \\
Shape actuator & Nitinol (NiTi) & Nitinol (NiTi) \\
Transition temperature ($^\circ$C) & 37 & 37 \\
Bonding adhesive & PEO (water-soluble) & PEO (water-soluble) \\
\bottomrule
\end{tabular}
\end{table}

\subsection{Impedance Stability}

The 64-channel SCEA maintained impedance stability through the compression--expansion shape transformation cycle (Table~\ref{tab:impedance}).

\begin{table}[H]
\centering
\caption{Impedance characterization of the 64-channel SCEA before and after one compression/expansion cycle. Measured in vitro at 1~kHz in phosphate-buffered saline. Paired $t$-test across all 64 channels. Cohen's $d$ from paired differences.}
\label{tab:impedance}
\begin{tabular}{lccccc}
\toprule
\textbf{Metric} & \textbf{Pre-Transform} & \textbf{Post-Transform} & $\boldsymbol{t(63)}$ & $\boldsymbol{p}$ & \textbf{Cohen's $d$} \\
\midrule
Mean impedance (k$\Omega$) & $48.7 \pm 6.2$ & $49.3 \pm 5.8$ & 0.81 & 0.42 & 0.10 \\
CV (\%) & 12.7 & 11.8 & --- & --- & --- \\
Functional channels & 64/64 & 64/64 & --- & --- & --- \\
\bottomrule
\end{tabular}
\end{table}

\subsection{Longitudinal MRI Assessment}

Structural T2-weighted MRI showed no brain changes at any timepoint. DCE-MRI assessed blood--brain barrier (BBB) integrity and meningeal enhancement (Table~\ref{tab:mri}).

\begin{table}[H]
\centering
\caption{Longitudinal MRI assessment in rats after SCEA implantation. T2-weighted structural MRI and DCE-MRI for BBB assessment. Enhancement = gadolinium signal enhancement ratio vs.\ pre-injection baseline. $n = 6$ rats. Friedman test for meningeal enhancement across timepoints: $\chi^2(4) = 18.7$, $p < 0.001$.}
\label{tab:mri}
\begin{tabular}{lcccc}
\toprule
\textbf{Timepoint} & \textbf{Structural Change} & \textbf{Cortical BBB Breach} & \textbf{Meningeal Enhancement} & \textbf{Enhancement Ratio} \\
\midrule
1 week & None & None & Present & $1.42 \pm 0.18$ \\
2 weeks & None & None & Diminishing & $1.28 \pm 0.14$ \\
3 weeks & None & None & Minimal & $1.12 \pm 0.09$ \\
4 weeks & None & None & Resolved & $1.04 \pm 0.05$ \\
8 weeks & None & None & Absent & $1.01 \pm 0.03$ \\
\bottomrule
\end{tabular}
\end{table}

\subsection{Immunohistochemistry}

IHC confirmed minimal inflammatory response at 8 weeks, comparable to standard epidural implants (Table~\ref{tab:ihc}).

\begin{table}[H]
\centering
\caption{Immunohistochemistry quantification in rat cortex beneath SCEA vs.\ contralateral control. Fold change = SCEA side / contralateral side. Values near 1.0 indicate no excess inflammation. $n = 4$ rats per timepoint. Mann--Whitney $U$ test for SCEA vs.\ contralateral at each timepoint.}
\label{tab:ihc}
\begin{tabular}{llcccc}
\toprule
\textbf{Marker} & \textbf{Target} & \textbf{1 Week (fold)} & \textbf{4 Weeks (fold)} & \textbf{8 Weeks (fold)} & $\boldsymbol{p}$ \textbf{(8 wk)} \\
\midrule
GFAP & Astrocyte reactivity & $1.38 \pm 0.15$ & $1.14 \pm 0.09$ & $1.12 \pm 0.08$ & 0.14 \\
Iba1 & Microglial activation & $1.31 \pm 0.12$ & $1.11 \pm 0.07$ & $1.09 \pm 0.06$ & 0.21 \\
NeuN & Neuronal density & $0.98 \pm 0.04$ & $0.99 \pm 0.03$ & $1.00 \pm 0.03$ & 0.89 \\
\bottomrule
\end{tabular}
\end{table}

\subsection{Rat Seizure Recording}

The SCEA captured high-bandwidth micro-ECoG signals during 4-aminopyridine (4-AP)-induced seizures, including high-frequency oscillations up to 500~Hz (Table~\ref{tab:seizure}).

\begin{table}[H]
\centering
\caption{Signal features captured during 4-AP-induced seizures in rats. Epidural recording via 64-channel SCEA. SNR = signal-to-noise ratio relative to baseline. $n = 4$ rats, $n = 12$ seizure episodes total.}
\label{tab:seizure}
\begin{tabular}{lccc}
\toprule
\textbf{Signal Feature} & \textbf{Frequency Range (Hz)} & \textbf{SNR (dB)} & \textbf{Detection Rate (\%)} \\
\midrule
Low-frequency oscillations & 0.5--50 & $24.3 \pm 3.1$ & 100 \\
Ripples & 80--250 & $12.7 \pm 2.4$ & 91.7 \\
Fast ripples & 250--500 & $8.4 \pm 1.9$ & 75.0 \\
Spatiotemporal propagation & Broadband & N/A & 100 \\
\bottomrule
\end{tabular}
\end{table}

\subsection{Canine Subdural Implantation and Recording}

The 256-channel SCEA was successfully deployed subdurally in beagle dogs, recording cortical dynamics during emergence from anesthesia (Table~\ref{tab:canine}).

\begin{table}[H]
\centering
\caption{Canine subdural implantation results. 256-channel SCEA deployed through a small dura slit. $n = 3$ beagles. Brain state classification from spectral analysis of recorded ECoG.}
\label{tab:canine}
\begin{tabular}{lccc}
\toprule
\textbf{Brain State} & \textbf{Dominant Frequency (Hz)} & \textbf{Spatial Pattern} & \textbf{RMS Amplitude ($\mu$V)} \\
\midrule
Deep anesthesia & $0.5--2$ & Globally synchronized & $82 \pm 14$ \\
Light anesthesia & $4--8$ & Regional differences & $127 \pm 21$ \\
Emergence & $8--30$ & Spatiotemporally organized & $198 \pm 34$ \\
\bottomrule
\end{tabular}
\end{table}

\subsection{Comparison with Existing Techniques}

Table~\ref{tab:comparison} compares SCEA with established brain mapping methods.

\begin{table}[H]
\centering
\caption{Comparison of SCEA with existing brain mapping techniques. $\uparrow$ = higher is better; $\downarrow$ = lower is better. Bold indicates best performance in each category.}
\label{tab:comparison}
\begin{tabular}{lccccc}
\toprule
\textbf{Technique} & \textbf{Spatial Res.\ $\downarrow$} & \textbf{Temporal Res.\ $\downarrow$} & \textbf{Bandwidth $\uparrow$} & \textbf{Invasiveness $\downarrow$} \\
\midrule
Scalp EEG & cm & ms & 0.5--50 Hz & $\mathbf{0}$ (none) \\
MEG & mm & ms & $\leq$ 250 Hz & 0 (none) \\
fMRI & mm & s & Hemodynamic & 0 (none) \\
fNIR & mm--cm & s & Hemodynamic & 0 (none) \\
Standard ECoG & $\mu$m--mm & ms & $\leq$ 500 Hz & 3 (craniotomy) \\
\textbf{SCEA} & $\boldsymbol{\mu}$\textbf{m--mm} & \textbf{ms} & $\boldsymbol{\leq}$ \textbf{500 Hz} & \textbf{1 (slit)} \\
\bottomrule
\end{tabular}
\end{table}

\subsection{Material Properties Under Deformation}

Table~\ref{tab:materials} quantifies the critical material properties enabling shape transformation.

\begin{table}[H]
\centering
\caption{Material properties of SCEA components under mechanical deformation. Conductivity retention = post-deformation / pre-deformation conductivity. $n = 10$ test samples per material.}
\label{tab:materials}
\begin{tabular}{lccl}
\toprule
\textbf{Component} & \textbf{Conductivity Retention (\%)} & \textbf{Failure Strain (\%)} & \textbf{Role} \\
\midrule
CNT film alone & $97.2 \pm 1.4$ & $> 20$ & Preserves conductivity \\
Au layer alone (nm-thick) & $12.3 \pm 8.7$ & $\sim 2$ & Fails under strain \\
CNT/Au hybrid & $95.8 \pm 2.1$ & $> 15$ & Combined function \\
Parylene-C (4 $\mu$m) & N/A & $> 30$ & Flexible substrate \\
Nitinol actuator & N/A & $> 8$ (superelastic) & Shape recovery \\
\bottomrule
\end{tabular}
\end{table}

\section{Discussion}

We have demonstrated a shape-changing electrode array that enables minimally invasive large-scale intracranial brain mapping. The SCEA addresses the central challenge in clinical ECoG: the tension between array size and surgical invasiveness. By transforming from a narrow strip to a large sheet at body temperature, the SCEA achieves large-area cortical coverage through a surgical opening that is an order of magnitude smaller than conventional craniotomy.

Several design features are critical to the SCEA's function. The CNT/Au multilayer maintains conductivity under the extreme mechanical deformation of the compression--expansion cycle, while nanometer-thick gold alone would crack and fail. The 4-$\mu$m Parylene-C substrate provides the flexibility needed for conformal contact with the cortical surface while remaining robust during shape transformation. The nitinol actuator exploits the martensite-to-austenite phase transition at body temperature, enabling autonomous deployment without external power.

Biocompatibility results are encouraging. The absence of cortical BBB breach on DCE-MRI at any timepoint, combined with resolution of meningeal enhancement by 4 weeks and minimal glial scarring at 8 weeks, suggests that the minimally invasive insertion causes less tissue trauma than conventional craniotomy. The preservation of neuronal density (NeuN) beneath the implant further supports long-term biocompatibility.

The high-bandwidth recording capability is clinically significant. Detection of fast ripples (250--500~Hz) is increasingly recognized as an important biomarker for epileptogenic zones \citep{worrell2004high, jacobs2009high}. The ability to detect these signals through a minimally invasive procedure could substantially improve the risk-benefit ratio of invasive epilepsy monitoring.

Limitations include the current demonstration in only two animal models (rat epidural, canine subdural). Human translation will require addressing scalability, sterilization protocols, and long-term chronic recording stability beyond 8 weeks. The actuator withdrawal step adds complexity to the surgical procedure. Future work should explore chronic recording over months and integration with wireless data transmission.

\section{Methods}

\subsection{SCEA Fabrication}

The SCEA was fabricated on a silicon carrier wafer. A 4-$\mu$m Parylene-C film was deposited by CVD, followed by photolithography and deposition of CNT/Au bilayer for recording sites and interconnects. SU8 insulation layers were patterned to define active electrode sites ($100 \times 100$~$\mu$m$^2$). CNT thin films (web-like SWCNT) were grown by CVD. Nitinol actuators were pre-shaped in the austenite phase and compressed at low temperature (martensite phase). PEO adhesive bonded the actuator to the array.

\subsection{In Vitro Characterization}

Impedance measurements were performed in PBS at 1~kHz using an electrochemical impedance analyzer. Measurements were taken before and after one full compression--expansion cycle.

\subsection{Rat Experiments}

Sprague-Dawley rats ($n = 6$) underwent epidural SCEA implantation through a small skull slit. Longitudinal T2-weighted and DCE-MRI were acquired at 1, 2, 3, 4, and 8 weeks post-implantation. IHC (GFAP, Iba1, NeuN) was performed at 1, 4, and 8 weeks. Seizures were induced by topical 4-AP application through a separate cranial hole.

\subsection{Canine Experiments}

Beagle dogs ($n = 3$) underwent subdural SCEA implantation through a small dura slit under general anesthesia. The 256-channel array was deployed at body temperature. ECoG was recorded during emergence from anesthesia. Post-operative MRI with nickel markers (Ni-SCEA) was used for array localization.

\subsection{Statistical Analysis}

Paired $t$-tests compared pre- and post-transformation impedance. Friedman tests assessed longitudinal MRI changes. Mann--Whitney $U$ tests compared IHC markers between SCEA and contralateral hemispheres. All tests used $\alpha = 0.05$.

\section{Conclusion}

The shape-changing electrode array represents a fundamental advance in minimally invasive neural interface technology. By combining ultrathin flexible electronics, mechanically robust CNT conductors, and temperature-responsive shape-memory actuation, the SCEA enables large-scale, high-bandwidth cortical mapping through surgical openings an order of magnitude smaller than conventional approaches. The demonstrated biocompatibility, recording performance, and scalability to 256 channels position this technology for future clinical translation in epilepsy monitoring, brain-computer interfaces, and fundamental neuroscience.

\bibliographystyle{plainnat}
\bibliography{references}

\end{document}
